% Generated draft: requires mathematical and source review.
\documentclass[11pt]{article}
\usepackage[margin=25mm]{geometry}
\usepackage{fontspec}
\setmainfont{Latin Modern Roman}
\usepackage{amsmath,amssymb,mathtools,unicode-math}
\setmathfont{Latin Modern Math}
\usepackage{enumitem,needspace,xurl,hyperref}
\setlength{\parindent}{0pt}
\setlength{\parskip}{5pt}
\setlength{\emergencystretch}{3em}
\newcommand{\N}{\mathbb{N}}
\newcommand{\Z}{\mathbb{Z}}
\newcommand{\Q}{\mathbb{Q}}
\newcommand{\R}{\mathbb{R}}
\newcommand{\C}{\mathbb{C}}
\newcommand{\F}{\mathbb{F}}
\newcommand{\A}{\mathbb{A}}
\newcommand{\PP}{\mathbb{P}}
\newcommand{\E}{\mathbb{E}}
\newcommand{\eps}{\varepsilon}
\newcommand{\dd}{\,\mathrm{d}}
\begin{document}
\section*{UnsolvedMath ID: 30005883}
\textbf{Source problem code:} catalogue entry
\subsection{Definitions and mathematical statement}
Let \(G\) be a finite simple group, let \(\operatorname{Irr}(G)\) denote its irreducible complex characters, and write \(g^G=\{x^{-1}gx:x\in G\}\) for the conjugacy class of \(g\in G\). For \(\chi\in\operatorname{Irr}(G)\), its character ratio at \(g\) is \(\chi(g)/\chi(1)\). The supplied record mentions character-ratio bounds but supplies no complete inequality involving these ratios: its original excerpt begins with the detached fragment ``for all \(\chi\in\operatorname{Irr}(G)\).'' No missing inequality can therefore be reconstructed from this record.

For a specified random-walk transition matrix \(P\) on \(G\) with uniform stationary distribution, define, for \(0<\delta<1\),
\[
 t_{\mathrm{mix}}(\delta)=\inf\left\{t\in\Z_{\geq0}:\max_{x\in G}\frac12\sum_{y\in G}\left|P^t(x,y)-\frac1{|G|}\right|\leq\delta\right\},
\]
with the infimum of the empty set equal to \(+\infty\). Under the directed Cayley-graph convention, \(\Gamma(G,C)\) has edges \(x\to xc\), where \(x\in G\) and \(c\in C\). Its directed diameter is the maximum shortest directed-path length, taken to be infinite if some ordered pair is not connected by a directed path. The source does not specify whether its Cayley graph is directed, whether inverse steps or holding probabilities are included, or which mixing threshold is intended.

The supplied Lubotzky conjecture, also attributed to Shalev, asks for the following comparison when \(G\) is simple of Lie type:
\[
 t_{\mathrm{mix}}\asymp\operatorname{diam}\Gamma(G,g^G).
\]
The excerpt additionally states, for \(G=\mathrm{Cl}_n(q)\),
\[
 t_{\mathrm{mix}}=\Omega\!\left(\frac{n}{\operatorname{supp}(g)}\right).
\]
Here the symbols \(\asymp\) and \(\Omega\) express, respectively, comparison above and below by positive multiplicative constants, and a lower bound by a positive multiplicative constant, in an intended asymptotic family. The record does not specify that family, the allowed dependence of these constants, the permitted elements \(g\), the meaning of \(n\), or the definition of \(\operatorname{supp}(g)\). In particular, a positive support and a generating conjugacy class are needed for the displayed finite comparison to have its usual meaning; these hypotheses are not supplied. The lower-bound direction \(\Omega\) is retained exactly as given. The source says that the diameter part was established in its reference [19], but does not provide the precise established theorem.

Let \(\alpha\) be a character of a finite-dimensional complex representation of \(G\). Write
\[
 \alpha=\sum_{\chi\in\operatorname{Irr}(G)}m_\chi\chi,
 \qquad
 \alpha^*=\sum_{m_\chi>0}\chi,
 \qquad
 \alpha^*(1)=\sum_{m_\chi>0}\chi(1).
\]
The character is faithful if the corresponding representation has trivial kernel. Under the usual directed convention, the McKay graph \(M(G,\alpha)\) has vertex set \(\operatorname{Irr}(G)\), with an edge \(\chi\to\psi\) whenever \(\psi\) occurs in the product character \(\alpha\chi\). The supplied Liebeck--Shalev--Tiep conjecture, attributed to reference [20], is written as
\[
 \operatorname{diam}M(G,\alpha)
 =\Omega\!\left(\frac{\log|G|}{\log\alpha^*(1)}\right)
 \qquad
 \text{for simple }G\text{ and faithful }\alpha.
\]
This is the literal lower-bound direction in the supplied excerpt. Its asymptotic range, constant dependence, and directed or undirected diameter convention are unspecified. The quotient requires \(\alpha^*(1)>1\). The record does not explicitly restrict simple groups to nonabelian simple groups, so this denominator condition cannot silently be assumed for every stated case. The supplied broad request for sharp bounds does not provide a separate upper-bound formula.

The excerpt also records claimed progress on Thompson's conjecture. For \(C\subseteq G\), put \(C^2=\{ab:a,b\in C\}\); the formulation given is the existence of a conjugacy class \(C\) such that \(G=C^2\). It reports that reference [3] proves this for all finite simple groups of Lie type over fields with at least nine elements. Its Theorem 4, attributed to [15] and [16], states that it holds for a simple classical group \(G=\mathrm{Cl}_n(q)\) of sufficiently large rank, subject to the additional conditions
\[
 G=\mathrm{P}\Omega_n^{\epsilon}(q),\quad 2\nmid q
 \quad\Longrightarrow\quad
 2\mid n\ \text{ and }\ \epsilon=(-1)^{n(q-1)/4}.
\]
The rank threshold is not supplied. The excerpt says this progress relies on its Theorem 3 and references [17] and [18], whose statements and bibliographic identities are absent. These are supplied contextual claims, not independently verified theorems in this draft.
\subsection{Short English statement}
The record asks how conjugacy-class random walks and McKay-graph diameters behave for finite simple groups. It supplies a mixing-time comparison and two bounds explicitly written with \(\Omega\), but omits several conventions and hypotheses needed for a unique precise formulation. It also reports progress on Thompson's conjecture. Its partially solved status is supplied dated curation, and this statement-only draft requires human mathematical review.
\subsection{Sources}
\begin{itemize}
\item Local UnsolvedMath registry: permanent ID \(30005883\), exact source problem code \texttt{catalogue entry}, title ``Mixing Times and Character Bounds for Simple Groups,'' in the Oberwolfach Reports Open Problems collection. The record gives proposed year 2024 and no named proposer.
\item Supplied report citation: \emph{Mini-Workshop: Growth and Expansion in Groups} (2024), Oberwolfach Reports, pp. 1005--1040, DOI 10.4171/owr/2024/17. \url{https://doi.org/10.4171/owr/2024/17}.
\item Supplied literature source, ``Uniform Character Bounds'': \url{https://annals.math.princeton.edu/2024/200-1/p01}.
\item Supplied literature source, ``probabilistic Waring results'': \url{https://annals.math.princeton.edu/2019/190-2/p03}.
\item Supplied literature source, ``OWR report'': \url{https://ems.press/content/serial-article-files/51362?nt=1}.
\item Literature assessment dated exactly 2026-08-22: the supplied curation labels the problem partially solved and reports that diameter and many Thompson and McKay consequences are proved, while the claimed mixing-time comparison and the full optimal character-bound program remain unsettled in general. This is supplied dated curation derived from the OWR material, not an independently verified literature assessment; the record explicitly requests independent verification of time-sensitive status and cited claims.
\item Attribution within the excerpt: the first conjecture is attributed to Lubotzky and also Shalev, with reference labels [23] and [24]; the McKay conjecture is attributed to Liebeck--Shalev--Tiep, with label [20]. Other cited labels are [3], [15], [16], [17], [18], and [19]. Their bibliographic entries are not supplied, so no identification with the linked literature sources is asserted.
\item The excerpt acknowledges support from the NSF, grant DMS-2200850, the Simons Foundation, and the Joshua Barlaz Chair in Mathematics. The acknowledging author's name is not supplied.
\end{itemize}
\end{document}
